% Clase 28 --- El truco de Young (§3.1).
% Sin la rendija previa cada punto de la fuente hace su propio patron, corrido
% respecto de los demas, y la suma se lava. La rendija fina selecciona un solo
% punto: todo lo que pasa por ella es coherente consigo mismo.
% Paneles APILADOS: cada uno es ancho y de lado a lado no entran en 16 cm.
\begin{center}
\begin{tikzpicture}[scale=0.9, transform shape]

% ---------------- (a) fuente extensa, sin rendija previa ----------------
\begin{scope}
  \foreach \yy in {-0.6,-0.3,0,0.3,0.6}{
    \filldraw[figamber] (-3.4,\yy) circle (1.6pt);
    \draw[figvecaux, figamber] (-3.25,\yy) -- (-0.15,{0.6+0.15*\yy});
    \draw[figvecaux, figamber] (-3.25,\yy) -- (-0.15,{-0.6+0.15*\yy});}
  \node[figlbl, figamber, anchor=east] at (-3.6,0) {vela};

  \draw[figline, line width=1.4pt] (0,-1.6) -- (0,-0.72);
  \draw[figline, line width=1.4pt] (0,-0.48) -- (0,0.48);
  \draw[figline, line width=1.4pt] (0,0.72) -- (0,1.6);

  \foreach \yy in {-1.3,-0.9,...,1.4}{
    \draw[figvecaux] (0.3,{0.3*\yy}) -- (2.9,\yy);}
  \fill[figgray!30] (3.0,-1.6) rectangle (3.4,1.6);
  \draw[figgray, line width=0.7pt] (3.0,-1.6) rectangle (3.4,1.6);
  \node[figlbl, anchor=west, align=left] at (3.6,0)
    {cada punto corre el patrón:\\ la pantalla queda pareja};

  \node[figlbl, anchor=north, align=center] at (0.2,-1.95)
    {\textbf{(a)} fuente extensa, sin $S_0$: incoherente};
\end{scope}

% ---------------- (b) con la rendija previa ----------------
\begin{scope}[yshift=-5.3cm]
  \foreach \yy in {-0.6,-0.3,0,0.3,0.6}{
    \filldraw[figamber] (-3.4,\yy) circle (1.6pt);
    \draw[figvecaux, figamber] (-3.25,\yy) -- (-1.75,0);}
  \node[figlbl, figamber, anchor=east] at (-3.6,0) {vela};

  \draw[figline, line width=1.4pt] (-1.65,-1.6) -- (-1.65,-0.12);
  \draw[figline, line width=1.4pt] (-1.65,0.12) -- (-1.65,1.6);
  \node[figlblaux, anchor=south] at (-1.65,1.68) {$S_0$};

  \draw[figvecres] (-1.6,0) -- (-0.15,0.6);
  \draw[figvecres] (-1.6,0) -- (-0.15,-0.6);

  \draw[figline, line width=1.4pt] (0,-1.6) -- (0,-0.72);
  \draw[figline, line width=1.4pt] (0,-0.48) -- (0,0.48);
  \draw[figline, line width=1.4pt] (0,0.72) -- (0,1.6);

  \foreach \yy in {-1.3,-0.9,...,1.4}{
    \draw[figvecaux] (0.3,{0.3*\yy}) -- (2.9,\yy);}
  \draw[figgray, line width=0.7pt] (3.0,-1.6) rectangle (3.4,1.6);
  \foreach \yy in {-1.35,-0.81,...,1.4}{
    \fill[figamber!80] (3.0,{\yy-0.1}) rectangle (3.4,{\yy+0.1});}
  \node[figlbl, anchor=west, align=left] at (3.6,0)
    {un solo punto emisor:\\ franjas nítidas};

  \node[figlbl, anchor=north, align=center] at (0.2,-1.95)
    {\textbf{(b)} el arreglo de Young (1801)};
\end{scope}

\end{tikzpicture}\\[0.5em]
{\small\itshape Young hizo el experimento en 1801, más de un siglo antes del
láser, y el problema que tuvo que resolver primero fue el de la
\textbf{coherencia}: la llama de una vela es una fuente extensa, cada pedacito
emite por su cuenta y con fase propia, de modo que cada uno produce su patrón de
franjas ligeramente corrido respecto del vecino. La suma de todos ellos es un
manchón parejo, como en (a). La solución fue interponer una \textbf{rendija muy
fina} antes de la doble rendija: todo lo que la atraviesa proviene
esencialmente de un único punto y es, por lo tanto, coherente consigo mismo, así
que las dos rendijas quedan alimentadas por la misma señal. Hay un compromiso
inevitable: cuanto más fina es $S_0$, mejor la coherencia y menos luz pasa; una
rendija infinitamente fina daría un patrón perfecto y completamente invisible.
Un láser resuelve las dos cosas a la vez ---es intenso y coherente en una región
amplia--- y por eso hoy el experimento sale a la primera.}
\end{center}
